\chapter{Computer Ethics}

\section{Codes of Ethics}
Codes of ethics vary from one professional society to another, but they typically share common features in prescribing the responsibilities of engineers to the public, their employers and clients, and their fellow engineers.

\subsection{Key Principles of Software Engineer Ethics}
\begin{enumerate}[itemsep=1ex]
   \item \textbf{Public Interest:} Software engineers must prioritize public well-being by ensuring safe, ethical, and accessible software, disclosing risks, balancing interests, and contributing to societal good through responsible and collaborative efforts.
   \item \textbf{Client and Employer:} Software engineers must act in the best interests of their client and employer by maintaining competence, integrity, confidentiality, legal compliance, authorized use of resources, ethical reporting, and avoiding conflicts of interest.
   \item \textbf{Product:} Software engineers must ensure their products meet the highest standards by prioritizing quality, addressing ethical and practical concerns, following professional guidelines, and ensuring thorough testing and documentation to meet user needs.
   \item \textbf{Judgment:} Software engineers must uphold integrity and independent judgment by tempering technical decisions with human values, endorsing only work within their competence, remaining objective, avoiding deceptive practices, and disclosing any conflicts of interest.
   \item \textbf{Management:} Software engineers must ensure effective project management by promoting quality, adhering to standards, assigning tasks based on competence, and providing realistic estimates, protecting confidentiality, respecting intellectual property, and ensuring transparency and fairness in all practices.
   \item \textbf{Profession:} Software engineers must advance the profession's integrity by adhering to ethical standards, promoting public knowledge, supporting colleagues, avoiding conflicts of interest, and responsibly addressing errors or violations of the code.
   \item \textbf{Colleague:} Software engineers must treat colleagues fairly and supportively by encouraging ethical adherence, assisting professional development, crediting others' work, providing objective reviews, addressing concerns, sharing best practices.
   \item \textbf{Self:} Software engineers must commit to lifelong learning, improving their skills in creating safe, reliable, and quality software, understanding standards and laws, and promoting ethics while avoiding unfair treatment, influencing unethical actions.
\end{enumerate}

\subsection{How do we know what is ethical?}

\textbf{Verify:} 
\Xmark{?}
Is it sensible?
\Xmark{?}
Is it desirable?
\Xmark{?}
Will it bring greatest good to people?
\Xmark{?}
Will it violates others basic rights?

\textbf{Judgment:} Use of our own senses to judge the action.


\section{Environmental Ethics}

\subsection{Green Software}
Green software, also known as sustainable software, is software designed and developed to minimize energy consumption and environmental impact throughout its lifecycle, from development to deployment and usage.

\subsection{Sustainable Software Engineering}
Sustainable Software Engineering focuses on creating software that is efficient and responsible across five key perspectives: environmental, social, individual, economic, and technical sustainability, ensuring software contributes positively to society, economy, and the environment.